\newpage
\begin{center}
  \textbf{\large ЗАКЛЮЧЕНИЕ}
\end{center}
\refstepcounter{chapter}
\addcontentsline{toc}{chapter}{ЗАКЛЮЧЕНИЕ}

В настоящей работе исследован метод физико-информированных нейронных сетей
применительно к задаче всплытия одиночного газового пузырька в жидкости. Задача
рассматривалась в двумерной постановке, соответствующей количественному
benchmark из работы Hysing et al. \cite{Hysing2009}. Эталонные данные получены
с помощью CFD-кода Basilisk с адаптивным измельчением сетки.

В работе реализованы и сопоставлены две архитектуры. Простая PINN, принимающая
на вход пространственно-временные координаты, демонстрирует хорошее качество
на обучающем режиме (\(L_2\) по \(\alpha\) около 7\%), но неспособна
распространить обученное решение на другие начальные радиусы: ошибка на
тестовом радиусе возрастает до 55\%. Параметрическая PINN, в которой начальный
радиус пузырька подается в сеть в качестве дополнительного входа, позволяет
одной модели описывать семейство задач с разными начальными условиями. При
тестировании на невиданных в обучении радиусах \(R=0.20\) и \(R=0.30\)
параметрическая модель показывает ошибки порядка 6\% по объемной доле и 6--11\%
по компонентам скорости, что на порядок лучше, чем у непараметрического подхода.

В ходе исследования проведена серия экспериментов:
сравнение трёх размеров архитектуры (light, medium, hard), анализ влияния числа
стадий обучения, а также аблация числа обучающих точек. Наиболее ёмкая
архитектура (8 слоёв по 350 нейронов) стабильно показывает лучшее качество по
всем физическим полям, кроме давления, где все варианты дают сопоставимые
результаты. Увеличение числа стадий оптимизации и числа PDE-точек каждое по
отдельности даёт умеренное улучшение метрик.

Основной вывод работы состоит в том, что параметрическая PINN представляет
собой жизнеспособный суррогат для быстрой оценки полей двухфазного течения при
изменении начального радиуса пузырька внутри изученного диапазона. Такой подход
может рассматриваться как элемент цифрового двойника установки: CFD обеспечивает
точные данные для обучения, а нейросетевая модель обеспечивает предсказание
на новых режимах без повторного решения нестационарной гидродинамической задачи.
